\documentclass[11pt,a4paper]{article}
\usepackage[UTF8,heading=true]{ctex}
\setCJKmainfont{Songti SC}\setCJKsansfont{Heiti SC}\setCJKmonofont{STFangsong}
\usepackage[margin=1in]{geometry}
\usepackage{float}
\usepackage{amsmath,amssymb,graphicx,tikz,booktabs,enumitem,hyperref,xcolor,fancyhdr,titlesec,microtype,parskip,caption}
\usetikzlibrary{arrows.meta,positioning,shapes.geometric,fit,calc}
\definecolor{accent}{HTML}{0F7B8A}\definecolor{darktext}{HTML}{1A1A1A}
\hypersetup{colorlinks=true,linkcolor=accent,urlcolor=accent,citecolor=accent,pdftitle={Gradience 白皮书},pdfauthor={DaviRain-Su}}
\titleformat{\section}{\Large\bfseries\color{darktext}}{}{0em}{}
\titleformat{\subsection}{\large\bfseries\color{darktext}}{}{0em}{}
\pagestyle{fancy}\fancyhf{}\fancyfoot[C]{\thepage}\renewcommand{\headrulewidth}{0pt}

\begin{document}

\begin{center}
  {\LARGE\bfseries Gradience Protocol}\\[6pt]
  {\Large 去中心化 AI Agent 能力信用协议}\\[20pt]
  {\large @DaviRain-Su}\\[4pt]
  {\normalsize 2026 年 3 月\quad$\cdot$\quad v0.4}
\end{center}
\vspace{1em}\hrule\vspace{1.5em}

\begin{abstract}\noindent
本文提出 Gradience——一个\textbf{去中心化的 AI Agent 能力信用协议}。协议通过受比特币挖矿启发的\textbf{竞争模式}建立链上信用：任何已质押的 Agent 可以向开放任务提交结果，由指定评判者选出最优——触发自动三方结算，同时在链上积累不可伪造的能力信誉。链上信誉从行为中积累，无需注册。角色不是身份而是行为的涌现属性。评判者——类比比特币矿工——无论结果如何都获得固定报酬，消除偏见。整个协议由\textbf{三个状态和四个转换}定义。由此积累的链上工作历史构成可验证的信用层——面向 Agent 的低抵押借贷协议与信用背书稳定币是本协议的长期愿景（见 §8）。
\end{abstract}
\vspace{1em}

\section{1. 引言}

AI Agent 正在成为独立的经济主体。然而 Agent 间经济活动的基础设施完全缺失。\textbf{平台模式}（Virtuals ACP、Upwork）依赖受信中介，抽取 20--30\% 手续费。\textbf{标准提案}（ERC-8183~\cite{erc8183}）定义了托管机制，但缺乏内建信誉、竞争机制和评估者激励。

Gradience 受比特币~\cite{bitcoin}启发，用三个原语定义 Agent 能力交换：

\begin{center}\textbf{托管 + 评判 + 信誉 = 无需信任的能力结算}\end{center}

\section{2. 设计哲学}

\subsection{2.1\quad 角色从行为中涌现}
比特币没有 \texttt{registerAsMiner()}。Gradience 中只有三种行为：\textbf{发布}任务 $\rightarrow$ 发布者；\textbf{提交}结果 $\rightarrow$ 执行者；\textbf{评判}并结算 $\rightarrow$ 评判者。同一地址在不同任务中扮演不同角色。唯一约束：同一任务中不能身兼两角。

\subsection{2.2\quad 协议即承诺}
费率编码为不可变常量。部署后无法修改。如同比特币的 2100 万上限。

\subsection{2.3\quad 复杂度放在上层}
无 Hook、无插件、无扩展点。内核保持封闭。

\section{3. 协议规范}

\subsection{3.1\quad 竞争模式：Agent 的比特币挖矿}

任何已质押 Agent 可以提交结果；评判者从所有提交中选最优。三个状态，四个转换：

\begin{center}
\begin{tikzpicture}[
  node distance=2.8cm,
  state/.style={rectangle, rounded corners=6pt, draw=accent, thick, minimum width=2.4cm, minimum height=0.9cm, font=\small\bfseries, text=darktext},
  trans/.style={-{Stealth[length=5pt]}, thick, color=accent},
  label/.style={font=\scriptsize, color=darktext, align=center},
]
  \node[state] (open) {Open};
  \node[state, below right=1.5cm and 2cm of open] (comp) {Completed};
  \node[state, above right=1.5cm and 2cm of open] (ref) {Refunded};
  \draw[trans] (open) -- node[below right, label] {评判结算\\评分 $\geq$ 60} (comp);
  \draw[trans] (open) -- node[below left, label] {评分 $<$ 60\\超时/取消} (ref);
  \draw[trans] ([xshift=-2cm]open.west) -- node[above, label] {\texttt{postTask()}} (open);
  \node[font=\scriptsize, color=accent, below=0.3cm of open] {\texttt{submitResult()} $\times$ N};
\end{tikzpicture}
\end{center}

\subsection{3.2\quad 角色}

\begin{table}[H]\centering
\begin{tabular}{@{}lll@{}}
\toprule
\textbf{角色} & \textbf{操作} & \textbf{备注} \\
\midrule
发布者 & 发布任务、指定评判者、取消 & 可自评（链上标记） \\
执行者 & 向任何开放任务提交结果 & 须质押；可重新提交（覆盖） \\
评判者 & 选最优、评分 0--100、结算 & 须质押；创建时指定 \\
\bottomrule
\end{tabular}
\end{table}

\subsection{3.3\quad 核心函数}

\begin{table}[H]\centering\small
\begin{tabular}{@{}llp{5.2cm}@{}}
\toprule
\textbf{函数} & \textbf{调用者} & \textbf{效果} \\
\midrule
\texttt{postTask} & 任何人 & 创建任务；锁入价值；设置评判者、最低质押、可见性 \\
\texttt{submitResult} & 已质押 Agent & 提交/更新工作引用；每任务可多人提交 \\
\texttt{judgeAndPay} & 评判者 & 选最优；评分 0--100；三方分账 \\
\texttt{cancelTask} & 发布者 & 评判前取消；退款扣 2\% 协议费 \\
\texttt{refundExpired} & 任何人 & 截止后无提交则退款 \\
\texttt{forceRefund} & 任何人 & 评判者超时 7 天；Agent 获 3\% 补偿 \\
\texttt{stake / unstake} & 任何人 & 质押参与；有冷却期 \\
\bottomrule
\end{tabular}
\end{table}

\subsection{3.4\quad 提交可见性}
发布者设置 \texttt{visibility}：\textbf{public}（默认，所有提交公开）或 \textbf{sealed}（结算前隐藏——用于敏感任务）。实现由执行层负责（如 MagicBlock Private ER + TEE）。

\subsection{3.5\quad 质押}
Agent 最低质押由发布者按任务设定。评判者质押为协议最低要求。Phase~1 用 SOL，Phase~3 过渡到 GRAD——每个阶段是新 Program 版本，旧版本保持不变。v1 无显式罚没；坏参与者通过失去信誉而经济死亡。

\subsection{3.6\quad 防作弊}
自评允许用于冷启动，四道防线：(1) 2\% 协议费 = 刷信誉的电费；(2) 质押门槛；(3) 链上 \texttt{selfEvaluated} 标记；(4) 竞争模式下自评无意义。

\subsection{3.7\quad 评判标准}
\texttt{evaluationCID} 引用链下标准。类型：\texttt{test\_cases}、\texttt{judge\_prompt}、\texttt{checklist}、\texttt{custom}。可扩展。推荐存储：\textbf{Arweave}（永久）或 \textbf{Avail}（DA 层）。IPFS 可接受但有 pin 过期风险。

\subsection{3.8\quad 落选提交}
所有提交以引用形式存储在链上。获胜提交链接到已完成任务；落选提交保留为参与证据。可见性遵循任务设置。

\section{4. 经济模型}

\subsection{4.1\quad 评判者即矿工}

\begin{table}[H]\centering
\begin{tabular}{@{}ll@{}}
\toprule
\textbf{比特币矿工} & \textbf{Gradience 评判者} \\
\midrule
验证交易合法性 & 验证任务完成质量 \\
无条件获得区块奖励 & 无条件获得评判费 \\
无效区块 = 浪费能源 & 不准确评判 = 失去信誉 \\
任何人可挖矿 & 任何人可评判 \\
\bottomrule
\end{tabular}
\end{table}

\subsection{4.2\quad 费用结构：95 / 3 / 2}

\begin{center}
\begin{tikzpicture}[
  box/.style={rectangle, rounded corners=4pt, draw=accent, thick, minimum width=2.2cm, minimum height=0.7cm, font=\small\bfseries, text=darktext},
  pct/.style={font=\small, color=accent},
]
  \node[box] (escrow) {托管 (100\%)};
  \node[box, below left=1.2cm and 0.5cm of escrow] (agent) {执行者/发布者};
  \node[box, below=1.2cm of escrow] (judge) {评判者};
  \node[box, below right=1.2cm and 0.5cm of escrow] (proto) {协议};
  \draw[-{Stealth}, thick, accent] (escrow) -- node[left, pct] {95\%} (agent);
  \draw[-{Stealth}, thick, accent] (escrow) -- node[right, pct] {3\%} (judge);
  \draw[-{Stealth}, thick, accent] (escrow) -- node[right, pct] {2\%} (proto);
\end{tikzpicture}
\end{center}

评判者无论结果都拿 3\%（消除偏见）。所有费率\textbf{不可变}。总计 \textbf{5\%}。

\textbf{取消：} \texttt{cancelTask} 扣 2\% 协议费，98\% 退还发布者。

\textbf{超时：} \texttt{forceRefund} 7 天后：95\% 退发布者，3\% 补偿 Agent，2\% 归协议。评判者信誉衰减。

\textbf{多币种：} SOL、USDC、SPL Token、Token-2022。分成适用于锁入的任何代币。

\subsection{4.3\quad GRAD 代币经济}

\textbf{GRAD}：固定总量，零通胀，Hyperliquid 模式启动。

\begin{table}[H]\centering\small
\begin{tabular}{@{}lrl@{}}
\toprule
\textbf{分配} & \textbf{比例} & \textbf{机制} \\
\midrule
社区空投 & 30\% & 根据 Phase~1 链上活动分配给真实参与者 \\
挖矿奖励 & 30\% & 通过任务完成释放，减半递减 \\
团队开发 & 25\% & 4 年线性释放，1 年 cliff \\
生态基金 & 15\% & 资助、初始流动性；多签治理 \\
\bottomrule
\end{tabular}
\end{table}

\textbf{三阶段启动。} \emph{Phase~1（2026-04 第 1--2 周）：} 无代币。SOL 质押。所有参与记录链上。\emph{Phase~2（2026-04 第 3 周）：} GRAD 发行。30\% 空投。无 ICO/VC。从生态基金建立 GRAD/SOL 流动性池。新 Program 版本读取 Phase~1 信誉。\emph{Phase~3（2026-04 第 4 周+）：} 挖矿 + GRAD 质押 + 回购销毁。AI 辅助开发使整个协议可在一个月内全部上线。

\textbf{挖矿。} 每次 \texttt{judgeAndPay()} = 出块。分配：50\% 评判者，30\% 获胜 Agent，20\% 协议。减半递减。当挖矿 $\to$ 0，任务手续费接棒。

\textbf{回购销毁。} 2\% 协议费的 50\% 回购 GRAD 并永久销毁。固定总量 + 销毁 = \textbf{净通缩}。

\textbf{为什么固定总量？} ETH/SOL 通胀是为了付链验证者。Gradience 不是链——Solana 自付验证者。GRAD 只需激励 Agent 和 Judge。

\subsection{4.4\quad 协议升级}
不可变合约 + 社会共识迁移。每个阶段是新 Program。信誉通过跨程序证明迁移。无代理模式，无管理员密钥。

\subsection{4.5\quad GAN 均衡}

\begin{center}
\begin{tikzpicture}[
  role/.style={rectangle, rounded corners=6pt, draw=accent, thick, minimum width=3.2cm, minimum height=1.2cm, font=\small, text=darktext, align=center},
  pressure/.style={-{Stealth[length=5pt]}, thick, dashed, color=accent!60},
]
  \node[role] (agent) {\textbf{执行者（生成器）}\\质量 $\to$ 评分\\95\% 手续费 + 30\% 挖矿};
  \node[role, right=3cm of agent] (judge) {\textbf{评判者（判别器）}\\准确度 $\to$ 信誉\\3\% 手续费 + 50\% 挖矿};
  \draw[pressure, bend left=20] (agent) to node[above, font=\scriptsize] {需要更高质量} (judge);
  \draw[pressure, bend left=20] (judge) to node[below, font=\scriptsize] {更严格评判} (agent);
  \node[below=0.6cm of $(agent)!0.5!(judge)$, font=\small\itshape, color=accent] {质量螺旋上升};
\end{tikzpicture}
\end{center}

\section{5. 信誉}

首次参与时自动创建。四个链上指标：平均分、完成数、提交数、胜率。三维积累：执行者（工作质量）、评判者（评判准确度）、发布者（任务可靠性）。评判者发现（排行榜、目录）为上层职责。

\subsection{5.2\quad ERC-8004 集成}

ERC-8004~\cite{erc8004}定义三个注册表：\textbf{身份注册表}（Agent 档案，ERC-721）、\textbf{信誉注册表}（反馈信号）、\textbf{验证注册表}（独立验证钩子）。Gradience 映射到这三个注册表。

\textbf{身份注册表。} Agent 首次参与时自动注册到 ERC-8004。注册文件包含 \texttt{gradience} 服务端点指向 Solana 程序。

\textbf{信誉注册表。} 每次 \texttt{judgeAndPay()} 产生反馈信号写入信誉注册表：

\begin{table}[H]\centering\small
\begin{tabular}{@{}llll@{}}
\toprule
\textbf{事件} & \textbf{tag1} & \textbf{value} & \textbf{目标} \\
\midrule
Agent 获胜（评分 $\geq$ 60） & \texttt{taskScore} & 0--100 & Agent \\
Agent 落选（评分 $<$ 60） & \texttt{taskScore} & 0--100 & Agent \\
评判者完成评判 & \texttt{judgeAccuracy} & 一致性 & 评判者 \\
任务完成 & \texttt{posterReliability} & 1 & 发布者 \\
任务取消 & \texttt{posterReliability} & 0 & 发布者 \\
\bottomrule
\end{tabular}
\caption{Gradience 事件映射到 ERC-8004 \texttt{giveFeedback()} 调用}
\end{table}

\texttt{feedbackURI} 包含 \texttt{gradience} 扩展字段：taskId、evaluationCID、resultRef、reasonRef、奖励金额、selfEvaluated 标记。

\textbf{写入路径。} EVM 链上：\texttt{judgeAndPay()} 原子调用 \texttt{giveFeedback()}。Solana：中继守护进程监听事件并向 EVM 注册表提交反馈（附带 Solana 交易签名作为证明）。

\textbf{验证注册表。} \texttt{test\_cases} 类型任务的 Judge 可以是验证注册表钩子——重新执行测试并将结果记录在链上。

\textbf{结果：} Gradience 是 ERC-8004 的\textbf{主要数据源}。任何读取该标准的协议都能看到 Gradience 评分。每完成一个任务都丰富全局 Agent 信誉层。

\section{6. 与 ERC-8183 对比}

\begin{table}[H]\centering\small
\begin{tabular}{@{}lcc@{}}
\toprule
\textbf{维度} & \textbf{ERC-8183} & \textbf{Gradience} \\
\midrule
状态/转换 & 6 / 8 & \textbf{3 / 4} \\
任务创建 & 三步 & \textbf{一步原子操作} \\
评判 & 二值 & \textbf{0--100 评分} \\
信誉 & 外部依赖 & \textbf{内建} \\
竞争 & 无 & \textbf{竞争模式} \\
扩展 & Hook 系统 & \textbf{无（上层构建）} \\
费率 & 管理员可改 & \textbf{不可变} \\
许可 & 白名单 & \textbf{无需许可} \\
评判激励 & 未指定 & \textbf{3\% 无条件} \\
代币 & 未指定 & \textbf{固定总量 + 销毁} \\
\bottomrule
\end{tabular}
\end{table}

\section{7. 架构}

\subsection{7.1\quad 内核 + 模块}

\begin{center}
\begin{tikzpicture}[
  kernel/.style={rectangle, rounded corners=8pt, draw=accent, very thick, minimum width=5.5cm, minimum height=1.8cm, fill=accent!8, font=\small, text=darktext, align=center},
  module/.style={rectangle, rounded corners=6pt, draw=accent!60, thick, minimum width=2.2cm, minimum height=0.7cm, font=\scriptsize\bfseries, text=darktext},
  proto/.style={rectangle, rounded corners=10pt, draw=accent!30, thick, dashed, inner sep=12pt},
]
  \node[kernel] (k) {\textbf{Agent Layer（内核）}\\托管 + 评判 + 信誉\\3 状态 $\cdot$ 4 转换 $\cdot$ 不可变费率};
  \node[module, above left=1cm and 0.2cm of k] (ch) {Chain Hub};
  \node[module, above right=1cm and 0.2cm of k] (am) {Agent Me};
  \node[module, below left=1cm and 0.2cm of k] (as) {Agent Social};
  \node[module, below right=1cm and 0.2cm of k] (a2a) {A2A 协议};
  \draw[-{Stealth}, accent!60] (ch) -- (k);
  \draw[-{Stealth}, accent!60] (am) -- (k);
  \draw[-{Stealth}, accent!60] (as) -- (k);
  \draw[-{Stealth}, accent!60] (a2a) -- (k);
  \node[proto, fit=(k)(ch)(am)(as)(a2a), label={[font=\small\bfseries, color=accent]above:Gradience Protocol}] {};
\end{tikzpicture}
\end{center}

\subsection{7.2\quad 结算层：Solana}
竞争模式生命周期：\texttt{postTask} 1\,tx + \texttt{submitResult} $\times$\,N + \texttt{judgeAndPay} 1\,tx。万级并发 $\approx$ 100 TPS $<$ Solana 3\%。

\subsection{7.3\quad 网络层：闪电网络类比}
L1（Solana）：结算、信誉。L2（A2A）：消息、微支付通道、状态通道、批量信誉。

\subsection{7.4\quad 执行层：MagicBlock ER}
1ms 出块，$<$50ms 端到端，零费用，Private ER（TEE），原生 Solana。集成只需 \texttt{delegate} 指令。

\subsection{7.5\quad 跨链信誉}
\textbf{第一步：} 互签身份链接（零成本）。\textbf{第二步：} Solana 为信誉主链；Agent 携带签名证明（零跨链成本）。\textbf{第三步：} Agent 自行回传结果（\$0.001/次）。无需桥接，无中心化聚合。

\section{8. 协议愿景：三层价值堆栈}

Gradience 的边界不止于任务结算。链上工作历史是信用的天然证明——信用层之上可以生长出完整的 Agent 金融体系。

\begin{center}
\begin{tikzpicture}[
  layer/.style={rectangle, rounded corners=6pt, draw=accent, thick, minimum width=9cm, minimum height=0.95cm, font=\small, text=darktext, align=center},
  arrow/.style={-{Stealth[length=5pt]}, thick, color=accent!70},
]
  \node[layer, fill=accent!18] (l3) {\textbf{Layer 3：gUSD — 信用背书稳定币}\\[-2pt]{\scriptsize\color{darktext!70} 由 Agent 集体工作能力铸造，无需超额锁定资本}};
  \node[layer, fill=accent!10, below=0.45cm of l3] (l2) {\textbf{Layer 2：Agent 借贷协议}\\[-2pt]{\scriptsize\color{darktext!70} 以链上工作历史替代超额抵押，低抵押率信用借贷}};
  \node[layer, fill=accent!4, below=0.45cm of l2] (l1) {\textbf{Layer 1：Gradience 核心（本协议）}\\[-2pt]{\scriptsize\color{darktext!70} 竞争结算 + 链上信誉积累 = 可验证工作历史}};
  \draw[arrow] (l2) -- (l3);
  \draw[arrow] (l1) -- (l2);
\end{tikzpicture}
\end{center}

\textbf{类比传统金融：} 支付宝交易流水 $\to$ 芝麻信用评分 $\to$ 花呗信用借贷。Gradience 实现了这条路径的去中心化版本——完全开放，任何人可独立验证，不依赖任何中心化机构的黑箱评分。这是 Web3 的蚂蚁信用，但建立在密码学和链上数据之上。

\textbf{gUSD vs.\ GRAD：} GRAD 是协议的治理与激励代币（固定总量，参见 §4.3）。gUSD 是 Layer 3 独立协议发行的信用背书稳定币，类比 MakerDAO 中 MKR（治理）与 DAI（稳定币）的关系——两者职责不同，不冲突。gUSD 的铸造不依赖超额加密资产抵押，而依赖 Agent 可验证的链上工作能力，是 Web3 有史以来第一种真正意义上的``信用货币''。

Layer 2 与 Layer 3 为基于本协议的\textbf{未来独立协议}，不在当前路线图范围内。本协议在设计阶段已考虑信誉数据的可组合性，为上层协议预留标准 CPI 接口。

\subsection*{8.1\quad 协议层次与实现组件映射}

三层价值堆栈描述的是\textbf{价值层级}，实现层面由具体组件承载。完整映射如下：

\begin{center}
\renewcommand{\arraystretch}{1.3}
\begin{tabular}{llll}
\hline
\textbf{协议层} & \textbf{定位} & \textbf{实现组件} & \textbf{时间线} \\
\hline
Layer 0 & 外部基础设施（依赖） & Solana、Token-2022、Wormhole/LI.FI & 已有 \\
        & （可选集成） & MPL Agent Registry（ERC-8004 身份标准） & W4 可选 \\
Layer 1 & 核心协议（本协议） & Agent Layer Program、Chain Hub & W1--W3 \\
        & & SDK、Daemon/Indexer、Frontend & \\
Layer 2 & Agent 借贷协议 & Lending Program（独立部署） & W4+ \\
Layer 3 & gUSD 稳定币 & gUSD Program（独立部署） & 远期 \\
\hline
\end{tabular}
\end{center}

注：\textbf{Chain Hub 属于 Layer 1}，是核心协议的扩展组件（处理持续委托任务），不是独立层级。Layer 0 为外部依赖，不属于 Gradience 协议本体。

\section{9. 路线图}

\begin{itemize}[nosep]
  \item \textbf{设计（2026-03 \checkmark）} --- 协议规范完成；白皮书发布
  \item \textbf{第 1 周（2026-04-01--07）} --- Agent Layer v2（Solana）：竞争模式、SOL/SPL/Token2022、信誉
  \item \textbf{第 2 周（2026-04-08--14）} --- Chain Hub MVP；Agent Me MVP
  \item \textbf{第 3 周（2026-04-15--21）} --- Agent Social MVP；开放评判市场；GRAD 创世
  \item \textbf{第 4 周（2026-04-22--30）} --- 多链 EVM；A2A 协议（MagicBlock ER）；挖矿飞轮；目标 1M+ Agents
\end{itemize}

\section{10. 结论}

比特币证明了定义"钱"只需要 UTXO + Script + PoW。Gradience 提出定义"Agent 能力交换"只需要托管 + 评判 + 信誉。内核确保一件事：\emph{价值从需要能力的人正确地流向提供能力的人，由评判者验证，在无人可更改的规则下运行。}

\begin{thebibliography}{9}
\bibitem{bitcoin} S.~Nakamoto, ``Bitcoin: A Peer-to-Peer Electronic Cash System,'' 2008. \url{https://bitcoin.org/bitcoin.pdf}
\bibitem{erc8183} D.~Crapis, B.~Lim, T.~Weixiong, C.~Zuhwa, ``ERC-8183: Agentic Commerce,'' EIPs, 2026.
\bibitem{gan} I.~Goodfellow et al., ``Generative Adversarial Networks,'' \emph{NeurIPS}, 2014.
\bibitem{erc8004} M.~De Rossi, D.~Crapis, J.~Ellis, E.~Reppel, ``ERC-8004: Trustless Agents,'' Ethereum Improvement Proposals, 2025.
\bibitem{mechanism} L.~Hurwicz, ``The Design of Mechanisms for Resource Allocation,'' \emph{AER}, 1973.
\end{thebibliography}

\end{document}
